\makeabstract
    {
    SnO\ensuremath{_2}-Coated AFM Tips for Nanoscale Electrical Characterization of Perovskite Solar Cells 
    }
    {
    Rowan Bixler\textsuperscript{1}, Joseph Brown\textsuperscript{2}, Dr. Jeffery Mativetsky\textsuperscript{2}
    }
    {
    \textsuperscript{1} Department of Physics, Amherst College, Amherst, MA\\
\textsuperscript{2} Department of Physics, SUNY Binghamton, Binghamton, NY
    }
    {
    Currently, commercial solar devices do not exceed 30\% power conversion efficiency (PCE). Perovskites present a valuable alternative to traditional silicon solar cells. The tunable bandgap, energy conversion efficiency, and quick energy payback of perovskites are inherent to their unique crystalline structure. Improving perovskite solar cell technology requires direct analysis of the perovskite layer at the Electron Transport Layer (ETL) interface. This interface is analyzed using an Atomic Force Microscope (AFM), which moves a sharp probe tip across a surface to map topography and local conductivity at the nanoscale. This research aims to improve perovskite analysis and effectively increase PCE, actualizing solar{\textquoteright}s potential in the renewable energy landscape.
    }
    {
    Various methods were explored for depositing SnO\ensuremath{_2}, a common perovskite ETL, onto microscopic AFM probe tips. Film thickness, stoichiometry, and electron mobility were optimized. Although not viable for the AFM probe itself, spin coating is a standardized technique that was used as a target for feasible tip deposition methods. Substrates were coated to establish parameters that were then applied to AFM tips. Tested methods included sputtering, where deposition time and oxygen concentration were optimized, and dip coating, where withdrawal rate and solution concentration were adjusted. Thin films were characterized using AFM and UV-vis spectroscopy. Additionally, bare tips on coated substrates were compared to coated tips on bare substrates.
    }
    {
    Analysis revealed that increasing oxygen concentration during sputtering yielded a bandgap that approached the spin coated thin film. Meanwhile, dip coated substrates exhibited bandgaps and electrical characteristics more closely aligned with spin coating. Electrical characterization of SnO\ensuremath{_2}-coated AFM tips showed that ideal substrate deposition conditions are not directly transferable to ideal deposition conditions for tips. This mismatch in electrical behavior is likely tip films are significantly thinner than substrates films. Ultimately, the sputtered tips exhibited more uniform coatings than the dip coated ones.
    }
    {
    Although dip coating naturally yields a bandgap and electrical characteristics closer to spin coating, sputtering has greater potential due to its tunability and also results in a compact layer that has not yet been achieved through dip coating. Future work will optimize sputtering for improved electrical properties, refine dip coating for sharp tip geometry, and compare coated versus bare tips on perovskite. These findings provide a foundation for advancing nanoscale characterization and realizing high-efficiency perovskite solar cells.
    }

    \newpage
    